\begin{abstract}
High-speed network receive paths face an acute tension between burst absorption and host last-level cache (LLC) efficiency. Direct cache injection (DDIO) keeps inbound packets in the LLC only while receive buffers stay within a few cache ways; traffic bursts that exceed that budget trigger Leaky DMA, evicting both packets and co-located application state to host DRAM. Existing solutions each leave a costly gap. Reactive rate control (HostCC) responds after cache misses begin. Host-capped rings (ShRing, rxBisect) remain bound by host memory conservation, dropping bursts or throttling transport rates. Recent dual-path on-NIC buffering (CEIO) absorbs bursts in NIC DRAM, yet splitting fast and slow paths requires per-flow steering rewrites, phase-exclusivity pauses, a two-producer software merge ring, and a host DMA-pull API.

We present \sys{}, a unified credit-gated receive architecture that collapses burst absorption and fast-path delivery onto a single FIFO push path. The host exposes only a strictly L3-pinned 256~KiB hot receive ring (4,096 cache-line cells), while an active NIC absorbs bursts in its own 48~MiB L3 and 6-channel DDR5 by default and pushes packets in order strictly against batched host credit returns. Built in Rust on an unmodified PCIe Gen5 server with a Marvell OCTEON~10 CN106 100GbE SmartNIC and verified across 1,089 zero-error hardware runs, \sys{} adds only 6~ns at zero backlog, maintains a flat 0.905~$\mu$s steady credit-refill hop across 128~MiB bursts, reduces host DRAM traffic by 98.9\% (100.0\% DDIO), and preserves co-located 32~MiB KV $p99$ latency at 41--42~ns (versus 153--259~ns for a 128~MiB host ring and 75.0--99.8\% drops for a passive 256~KiB ring). Integrated into a complete push/DMA datapath (DirectPush) over 100GbE, it saturates 100GbE RX on a single host core and improves dual-host RPC/KV throughput by up to 55\% while cutting open-loop $p99$ latency by 50--64\% over a tuned asynchronous descriptor-ring baseline.
\end{abstract}
